\newpage
%%%%%%%%%%%%%%%%%%%%%%%%%%%%%%%%%%%%%%%%%%%%%%%%%%%%%%%%%%%%%%%%
%%%%%%%%%%%%%%%%%%%%%%%%%%%%%%%%%%%%%%%%%%%%%%%%%%%%%%%%%%%%%%%%
%%%%%%%%%%%%%%%%%%%%%%%%%% Enunciado %%%%%%%%%%%%%%%%%%%%%%%%%%%

\begin{myblock}
\phantomsection\addcontentsline{toc}{section}{Parte B | LSTM}
\section*{Parte B | LSTM}

    Usando la misma base de la tarea 1, entrena un modelo por cada arquitectura: 
    CNN, RNN/LSTM/GRU, y un Transformer de Hugging Face con fine-tuning. Reporta
    accuracy y F1, incluye matriz de confusión y análisis de errores característicos.  

\end{myblock}

%%%%%%%%%%%%%%%%%%%%%%%%%%%%%%%%%%%%%%%%%%%%%%%%%%%%%%%%%%%%%%%%
%%%%%%%%%%%%%%%%%%%%%%%%%%%%%%%%%%%%%%%%%%%%%%%%%%%%%%%%%%%%%%%%

%%%%%%%%%%%%%%%%%%%%%%%%%%%%%%%%%%%%%%%%%%%%%%%%%%%%%%%%%%%%%%%%
%%%%%%%%%%%%%%%%%%%%%%%%%%%%%%%%%%%%%%%%%%%%%%%%%%%%%%%%%%%%%%%%

Si bien la arquitectrura interna de la celda LSTM es la misma tanto para 
generación de texto como para clasificación de texto, el objetivo final, y la
estructura del modelo, cambian. 


\subsection{Arquitectura LSTM}

En la tarea de generación, el modleo funciona como un Language Model. El objetivo
es predecir el siguiente token de la secuencia $y_i$, dado el contexto de los
tokens anteriores $(x_1, x_2,...,x_n)$. 

Como se implementó en el código y arquitectura de la parte A, la salida LSTM en 
cada paso a través del tiempo $h_t$ se pasa por una capa lineal de proyeción, seguida
de una Softmax sobre todo el vocabulario. La función de pérdida, como lo 
es la Entropía Cruzada, se calcula en cada token. De esa manera podemos hacer 
comparación de distribuciones de probabilidad predicah en el token real que
sigue en la secuencia. El modleo se entrena para minimizar dicha pérdida a lo 
largo de la secuencia.

\[
    \hat{y}_t = \mathrm{Softmax}(W_o h_t + b_o)
\]
\[
    \mathcal{L}_{\text{generación}} = \sum_{t=1}^{L} \mathrm{CrossEntropy}(\hat{y}_t, y_t)
\]

El objetivo es aprender la distribución de probabilidad $P(w_t | w_{1:t-1})$.

\subsubsection{Un Modelo de Representación Global}

Para clasificación de texto, no queremos predecir el siguiente token, queremos 
asignarle una etiqueta única a toda la secuencia o documento. Por lo tanto, no
nos interesan las salidad de la LSTM a lo largo del tiempo, sino una Representación
vectorial que resuma el contenido global de la secuencia completa. 

El modelo procesa la secuencia de entrada, para de esa manera obtener los estados
ocultos en cada paso, $H = [h_1, h_2,...,h_L]$. A partir de ahí, la diferencia
clave es la introducción de una capa de Pooling. 

\begin{itemize}
    \item \textbf{Agregación de Información (Pooling)}: en lugar de usar cada
    estado oculto $h_t$, se aplica una operación de pooling sobre todos los estados
    para obtener un único vector de tamaño fijo. Se pueden realizar diversos ajustes
    como la creación de un MeanPooling tal como vimos en RNN, pero el punto importante
    es que se toma un estado oculto final, que queda a modo de resumen de toda la
    secuencia que lo precede.
    \item \textbf{Proyección de Clases}: este vector resumen se pasa a la cabeza de la 
    clasificación, que es una capa lineal final que proyecta el vector al número de 
    clases deseadas, por ejemplo, 5 para la polaridad que tenemos de 1.0 a 5.0 estrellas.
\end{itemize}

\[
    \mathbf{v} = \mathrm{Pool}(h_1, h_2, \dots, h_L)
\]
\[
    \hat{\mathbf{y}}_{\text{clase}} = \mathrm{Softmax}(W_{\text{clf}} \mathbf{v} + b_{\text{clf}})
\]

La función de pérdida (`CrossEntropyLoss`) se calcula una sola vez por 
secuencia, comparando la predicción final con la etiqueta verdadera de
la secuencia.

\[
    \mathcal{L}_{\text{clasificación}} = \mathrm{CrossEntropy}(\hat{\mathbf{y}}_{\text{clase}}, \mathbf{y}_{\text{clase}})
\]

En resumen, mientras que el modelo generativo es un traductor 
secuencia-a-secuencia (de $N$ tokens a $N$ distribuciones de probabilidad), 
el modelo de clasificación es un codificador secuencia-a-vector, que condensa
toda la información de la secuencia en una única representación para tomar 
una decisión final. Esta distinción es la que justifica el cambio en la
arquitectura de salida y la forma en que se calcula la pérdida.

\subsection{Conclusiones}

La primera implementación de una BiLSTM de una sola capa con \emph{max-pooling} no solo superó a la RNN simple, sino que estableció un nuevo estado del arte para el problema, alcanzando un F1-macro en test superior a 0.51. Sin embargo, este éxito vino acompañado de un **sobreajuste extremo**, con métricas en el conjunto de entrenamiento que superaban el 90% mientras que en validación apenas rondaban el 50%. Este fenómeno motivó una serie de experimentos enfocados en dos frentes: aumentar la complejidad del modelo (más capas) y aplicar técnicas de regularización más robustas para mejorar la generalización.

\begin{table}[h!]
\centering
\caption{BiLSTM — Resumen de experimentos en MeIA (polaridad 1–5), transpuesto.}
\label{tab:lstm_meia_transposed}
\resizebox{\textwidth}{!}{%
\begin{tabular}{lcccc}
\toprule
\textbf{Parámetro} & \textbf{\texttt{...\_bilstm\_max\_lr1e-3}} & \texttt{...\_bilstm4\_max\_do20} & \texttt{...\_bilstm\_maxmean...} & \texttt{...\_bilstm\_max\_regstrong} \\
\midrule
\textbf{Capas LSTM}        & \textbf{1}          & 2                   & 2                & \textbf{1} \\
\textbf{Pooling}           & \textbf{max}        & max                 & max+mean         & \textbf{max} \\
\textbf{Regularización}    & Base                & Base                & Base             & Fuerte (D.O., WD, W.D.) \\
\textbf{Val F1-M}          & 0.495               & 0.493               & 0.489            & \textbf{0.505} \\
\textbf{Test Acc / F1-M}   & \textbf{0.514 / 0.515} & 0.480 / 0.461       & 0.474 / 0.470    & 0.484 / 0.477 \\
\bottomrule
\end{tabular}%
}
\end{table}

\subsection{Conclusiones}

\begin{itemize}
    \item El modelo baseline de BiLSTM de una capa se consolidó como el de mejor rendimiento de todas las arquitecturas recurrentes y convolucionales, 
    alcanzando un F1-macro de 0.515 en el conjunto de prueba.
    \item El sobreajuste fue el principal desafío. A pesar de los intentos por mitigarlo, ningún modelo regularizado superó al baseline en el conjunto de test,
    aunque sí lograron reducir la brecha entre las métricas de entrenamiento y validación.
    \item Aumentar la profundidad del modelo (usando 2 capas de BiLSTM) resultó contraproducente. Incrementó el tiempo de cómputo y la complejidad sin una
    mejora en la generalización, de hecho, el rendimiento en test disminuyó.
    \item La estrategia de regularización fuerte (mayor dropout, weight decay y word dropout) logró el mejor F1-macro en el conjunto de validación
    (0.505), indicando que el modelo estaba generalizando mejor a ese subconjunto de datos. Sin embargo, esta mejora no se tradujo al conjunto de test, donde el
    rendimiento fue inferior al del baseline.
    \item Se reitera la conclusión de secciones anteriores: el MaxPooling sigue siendo la estrategia de agregación más efectiva para este corpus,
    superando a la combinación \texttt{max+mean}.
\end{itemize}

La arquitectura BiLSTM demostró ser la más poderosa hasta el momento, con un modelo simple de una sola capa estableciendo un nuevo benchmark de rendimiento.
El análisis de los experimentos demuestra que el modelo más simple, a pesar de un sobreajuste masivo en el entrenamiento, encontró 
un mínimo en el espacio de parámetros que generalizó de forma más efectiva al conjunto de prueba final. Las intervenciones lógicas para mejorar el modelo, aumentar 
su capacidad con más capas o forzar la generalización con una regularización más estricta, paradójicamente degradaron su rendimiento en la evaluación final. 
Esto indica que las técnicas de regularización, si bien controlaron el sobreajuste, pudieron haber impedido al modelo aprender ciertos patrones cruciales 
presentes en el corpus, o que el conjunto de validación no fue un proxy perfecto para el de test. A pesar de todo, la BiLSTM se establece como el mejor modelo
de redes neuronales hasta el momenoto.

\subsection{Visualización de Modelos}

% --- Figura: Matrices de Confusión para LSTM ---
\begin{figure}[h!]
    \centering

    % --- PRIMERA FILA: DOS IMÁGENES ---
    \begin{subfigure}[b]{0.45\textwidth}
        \centering
        \includegraphics[width=\textwidth]{/home/cesar/project_b_t02nlp/models/lstm/exp_bilstm_max_lr1e-3/confusion_test.png} 
        \caption{\textbf{\texttt{...max\_lr1e-3}}}
        \label{fig:conf_lstm_exp1}
    \end{subfigure}
    \hfill % Espacio horizontal flexible entre las imágenes
    \begin{subfigure}[b]{0.45\textwidth}
        \centering
        \includegraphics[width=\textwidth]{/home/cesar/project_b_t02nlp/models/lstm/exp_bilstm4_max_do20/confusion_test.png}
        \caption{\texttt{...bilstm4\_max...}}
        \label{fig:conf_lstm_exp2}
    \end{subfigure}

    \vspace{1em} % Espacio vertical entre filas para mayor claridad

    % --- SEGUNDA FILA: UNA IMAGEN (CENTRADA) ---
    \begin{subfigure}[b]{0.45\textwidth}
        \centering
        \includegraphics[width=\textwidth]{/home/cesar/project_b_t02nlp/models/lstm/exp_bilstm_maxmean_ls000/confusion_test.png} 
        \caption{\texttt{...maxmean...}}
        \label{fig:conf_lstm_exp3}
    \end{subfigure}

    \vspace{1em} % Espacio vertical entre filas

    % --- TERCERA FILA: DOS IMÁGENES ---
    \begin{subfigure}[b]{0.45\textwidth}
        \centering
        \includegraphics[width=\textwidth]{/home/cesar/project_b_t02nlp/models/lstm/exp_bilstm_max_regstrong/confusion_test.png} 
        \caption{\texttt{...max\_regstrong}}
        \label{fig:conf_lstm_exp4}
    \end{subfigure}
    \hfill % Espacio horizontal flexible
    \begin{subfigure}[b]{0.45\textwidth}
        \centering
        \includegraphics[width=\textwidth]{/home/cesar/project_b_t02nlp/models/lstm/exp_bilstm_h192_do60_wd2e-4_wdrop005/confusion_test.png} 
        \caption{\texttt{...h192\_do60...}}
        \label{fig:conf_lstm_exp5}
    \end{subfigure}
    
    \caption{Matrices de confusión para los cinco experimentos de BiLSTM. El modelo baseline (primera imagen, arriba a la izquierda) muestra el mejor rendimiento.}
    \label{fig:confusion_matrices_lstm}
\end{figure}

%%%%%%%%%%%%%%%% CURVAS ROC %%%%%%%%%%%%%%%%%%
%%%%%%%%%%%%%%%%%%%%%%%%%%%%%%%%%%%%%%%%%%%%%%

% --- Figura: Curvas ROC OVR para LSTM ---
\begin{figure}[h!]
    \centering

    % --- PRIMERA FILA: DOS IMÁGENES ---
    \begin{subfigure}[b]{0.45\textwidth}
        \centering
        \includegraphics[width=\textwidth]{/home/cesar/project_b_t02nlp/models/lstm/exp_bilstm_max_lr1e-3/roc_test.png} 
        \caption{\textbf{\texttt{...max\_lr1e-3}}}
        \label{fig:roc_lstm_exp1}
    \end{subfigure}
    \hfill % Espacio horizontal flexible
    \begin{subfigure}[b]{0.45\textwidth}
        \centering
        \includegraphics[width=\textwidth]{/home/cesar/project_b_t02nlp/models/lstm/exp_bilstm4_max_do20/roc_test.png}
        \caption{\texttt{...bilstm4\_max...}}
        \label{fig:roc_lstm_exp2}
    \end{subfigure}

    \vspace{1em} % Espacio vertical para separar filas

    % --- SEGUNDA FILA: UNA IMAGEN (CENTRADA) ---
    \begin{subfigure}[b]{0.45\textwidth}
        \centering
        \includegraphics[width=\textwidth]{/home/cesar/project_b_t02nlp/models/lstm/exp_bilstm_maxmean_ls000/roc_test.png} 
        \caption{\texttt{...maxmean...}}
        \label{fig:roc_lstm_exp3}
    \end{subfigure}

    \vspace{1em} % Espacio vertical para separar filas

    % --- TERCERA FILA: DOS IMÁGENES ---
    \begin{subfigure}[b]{0.45\textwidth}
        \centering
        \includegraphics[width=\textwidth]{/home/cesar/project_b_t02nlp/models/lstm/exp_bilstm_max_regstrong/roc_test.png} 
        \caption{\texttt{...max\_regstrong}}
        \label{fig:roc_lstm_exp4}
    \end{subfigure}
    \hfill % Espacio horizontal flexible
    \begin{subfigure}[b]{0.45\textwidth}
        \centering
        \includegraphics[width=\textwidth]{/home/cesar/project_b_t02nlp/models/lstm/exp_bilstm_h192_do60_wd2e-4_wdrop005/roc_test.png} 
        \caption{\texttt{...h192\_do60...}}
        \label{fig:roc_lstm_exp5}
    \end{subfigure}
    
    \caption{Curvas ROC One-vs-Rest (OVR) para los cinco experimentos de BiLSTM. Se evalúa la capacidad de discriminación de cada clase.}
    \label{fig:roc_curves_lstm}
\end{figure}


\clearpage